% !TeX root = ../../Book5.tex
% 纲要标题：## 附录 D　导出范畴的几何应用与 Fourier–Mukai 变换
% 纲要位置：几何学计划纲要.md，第 6712 行。
% 正文状态：未撰写；本文件目前是框架。
% 证明范围与先修：以域上的光滑射影簇为主要范围，使用第18—22章的导出函子、上同调、基变换和对偶，第31章的完美复形及第33章的 Picard 语言。分别检查积分核的有界性、有限 Tor 维数、凝聚性、投影固有性与复合的基变换条件。椭圆曲线模型固定有原点的光滑射影曲线及归一化 Poincaré 丛，明确基域及对偶曲线；选定逆核与复合计算为代表模块，一般等价判据准确外引。P¹ 例外对用对角线消解说明正交性和生成性。不扩展为一般 Bridgeland 稳定性、DT 不变量或一般 motives 理论，也不以这些名称冒充已完成的范畴构造。

% 纲要内容开始
% 
% > 本附录证明范围与外部输入：以域上的光滑射影簇为主要范围，依赖第18—22章的导出函子、上同调、基变换与对偶，以及第31章完美复形和第33章 Picard 语言。分别核对积分核的有界性、有限 Tor 维数、凝聚性、投影固有性及复合的基变换条件。椭圆曲线模型固定原点及归一化 Poincaré 丛，明确基域、逆核和复合平移；一般等价判据准确外引。射影直线例外对以对角线消解完成正交性与生成性；Bridgeland、DT 和一般 motives 保持外部方向。
% 
% * **凝聚层的有界导出范畴**
%   * 复形、同伦与拟同构
%   * 局部化与有界导出范畴
%   * 平移、映射锥与三角
%   * 导出张量与导出 Hom
%   * 完美复形及光滑性条件
% 
% * **积分核与 Fourier–Mukai 变换**
%   * 乘积空间上的核与投影
%   * 有界性、凝聚性与推拉条件
%   * 核的卷积与变换复合
%   * 伴随变换与对偶
%   * 等价判据及其外部输入
% 
% * **椭圆曲线的 Poincaré 核模型**
%   * 对偶椭圆曲线与归一化 Poincaré 丛
%   * 点层与线丛的变换
%   * 上同调跳跃与基变换
%   * 逆核、复合与平移
%   * 模型等价的证明范围
% 
% * **P¹ 上的例外对与半正交分解**
%   * 例外对象与例外对
%   * 正交性与生成性
%   * 射影直线上的线丛对
%   * 对角线消解与分解
%   * 由模型到一般分解的边界
% 
% ---
% 
% 纲要内容结束

\chapter{导出范畴的几何应用与 Fourier–Mukai 变换}
\label{g5:app:D}
\BookChapterNumber{g5:app:D}

待撰写

\section{凝聚层的有界导出范畴}
\label{g5:sec:D01}
待撰写

\subsection{复形、同伦与拟同构}
\label{g5:subsec:D01:01}
待撰写

\subsection{局部化与有界导出范畴}
\label{g5:subsec:D01:02}
待撰写

\subsection{平移、映射锥与三角}
\label{g5:subsec:D01:03}
待撰写

\subsection{导出张量与导出 Hom}
\label{g5:subsec:D01:04}
待撰写

\subsection{完美复形及光滑性条件}
\label{g5:subsec:D01:05}
待撰写

\section{积分核与 Fourier–Mukai 变换}
\label{g5:sec:D02}
待撰写

\subsection{乘积空间上的核与投影}
\label{g5:subsec:D02:01}
待撰写

\subsection{有界性、凝聚性与推拉条件}
\label{g5:subsec:D02:02}
待撰写

\subsection{核的卷积与变换复合}
\label{g5:subsec:D02:03}
待撰写

\subsection{伴随变换与对偶}
\label{g5:subsec:D02:04}
待撰写

\subsection{等价判据及其外部输入}
\label{g5:subsec:D02:05}
待撰写

\section{椭圆曲线的 Poincaré 核模型}
\label{g5:sec:D03}
待撰写

\subsection{对偶椭圆曲线与归一化 Poincaré 丛}
\label{g5:subsec:D03:01}
待撰写

\subsection{点层与线丛的变换}
\label{g5:subsec:D03:02}
待撰写

\subsection{上同调跳跃与基变换}
\label{g5:subsec:D03:03}
待撰写

\subsection{逆核、复合与平移}
\label{g5:subsec:D03:04}
待撰写

\subsection{模型等价的证明范围}
\label{g5:subsec:D03:05}
待撰写

\section{\texorpdfstring{\(\mathbb P^1\) 上的例外对与半正交分解}{P¹ 上的例外对与半正交分解}}
\label{g5:sec:D04}
待撰写

\subsection{例外对象与例外对}
\label{g5:subsec:D04:01}
待撰写

\subsection{正交性与生成性}
\label{g5:subsec:D04:02}
待撰写

\subsection{射影直线上的线丛对}
\label{g5:subsec:D04:03}
待撰写

\subsection{对角线消解与分解}
\label{g5:subsec:D04:04}
待撰写

\subsection{由模型到一般分解的边界}
\label{g5:subsec:D04:05}
待撰写
